\section{A quartic in forty-five unrestricted integers}
\label{sec:quartic45}

The independent low-degree contribution reduces the inspected
fifty-integer quartic to \textbf{forty-five unrestricted integer unknowns},
while retaining degree four.  It uses 28 quadratic residuals instead of
29.  Its expanded representation has 407 symbolic terms, compared with
362 in the chosen parent; the improvement concerns variable count at
fixed degree, not monomial count or evaluation cost.

This construction keeps the complete package-34 coefficient tuple
\begin{center}\code{mixed-parity-unrounded-1064921},\end{center}
with
\(B=2^{1064921}\), \(\beta=B-1\), and ports
\(\beta,K_0,M_C,M_F\).  The last port is the already shifted
\(M_{F,\mathrm{native}}+B-1\).  It is a separate complete binding from
the tuple used for the principal twelve- and thirteen-variable equations.
The four raw payloads, their finite construction, the literal K215
transition table, and the relevant inherited proofs are included in
\texttt{quartic/upstream/}.  The canonical accepting witnesses for this
tuple have \(q=2^t\) with \(t\) odd.

\subsection{The two guard reductions}

Write
\[
 \operatorname{Pos}(u,v,t)=1+u^2+v^2+t(t+1).
\]
This expression is positive on all integer triples and represents every
positive integer.  For the converse, apply the three-square theorem to
\(4n-3\): its residue modulo four forces one odd square and two even
squares, yielding the displayed representation after division by four.

\begin{lemma}[A common-sign guard]
Suppose \(q>0\) and \(C=q-F-\alpha\).  Then
\[
 F\alpha>0\quad\hbox{and}\quad FC>0
 \quad\Longleftrightarrow\quad F,\alpha,C>0.
\]
\end{lemma}
\begin{proof}
The two products make the three nonzero integers have the same sign.
Their sum is \(q>0\), so that sign is positive.  The converse is immediate.
\end{proof}

\begin{lemma}[A dyadic guard]
For \(\beta\ge3\), the equation
\(\beta J+1=2v^2\), with \(J,v\in\mathbb Z\), implies \(J>0\).
Every canonical radix \(q=2^t\) with \(t\) odd satisfies this equation
with \(v=2^{(t-1)/2}\).
\end{lemma}
\begin{proof}
If \(J\le-1\), its left side is negative; if \(J=0\), it is one,
which cannot be twice an integer square.  The second assertion is a
direct substitution.
\end{proof}

\subsection{The complete quadratic system}

The 32 arithmetic coordinates are
\[
\begin{gathered}
J,F,\alpha,z,f,h,r,v_s,w,\tau,\eta,\zeta,y,\gamma,o,j,\\
N,E,P,c,a,X,D,C_2,L,F_2,V,g,I,B_{\rm aux},\Delta,u.
\end{gathered}
\]
Add one coordinate \(v_q\) and four triples
\(\boldsymbol e,\boldsymbol\xi,\boldsymbol b,\boldsymbol d\).
All 45 coordinates range over \(\mathbb Z\).  In the following formulas,
the letters
\[
\begin{gathered}
 q=\beta J+1,\quad C=q-F-\alpha,\quad
 k=\eta+\zeta,\quad H=4a+3,\\
 R=(q-F)(N-q)-CL+(M_C+qM_F)J
\end{gathered}
\]
are polynomial abbreviations, not additional coordinates.  Retain the
following 23 arithmetic equations:
\begin{alignat*}{2}
 L&=q^2-1,&\qquad N&=q(L+1),\\
 X&=wq,&a&=E+(2v_s+1)N,\\
 E&=X(a-E),&c&=k(a-E)+\eta,\\
 P&=E(c-\eta)+k,&\tau^2&=P(P-k)+1,\\
 D&=X+ac+\gamma H,&\Delta&=a^2+H,\\
 C_2&=c^2,&D^2&=\Delta C_2+1,\\
 k&=R+1+hE,&(K_0+w)C+q-F-z(q-1)&=1,\\
 F_2&=f^2,&V&=of-c,\\
 V-cj+k-1-hE&=0,&g&=y^2-V^2,\\
 I&=(2r+1)C_2,&B_{\rm aux}&=I^2,\\
 B_{\rm aux}&=\Delta(F_2-1),&y^2-B_{\rm aux}g&=1,\\
 w&=u^2.&&
\end{alignat*}
Add the five quadratic guards
\begin{align*}
 \eta&=\operatorname{Pos}(\boldsymbol e),&
 \zeta&=\operatorname{Pos}(\boldsymbol\xi),&
 q&=2v_q^2,\\
 F\alpha&=\operatorname{Pos}(\boldsymbol b),&
 FC&=\operatorname{Pos}(\boldsymbol d).
\end{align*}
Here \(\operatorname{Pos}(\boldsymbol e)\) means evaluation at its three
coordinates.  Every displayed residual is quadratic.  If they are
\(r_1,\ldots,r_{28}\), define the single equation
\begin{equation}
 \boxed{\displaystyle Q_{45}=\sum_{i=1}^{28}r_i^2=0.}
 \label{eq:quartic45}
\end{equation}
The coefficient of \(\tau^4\) is one, so the degree is exactly four.
The literal ordered residuals, variable list, and exact expansion are
\texttt{quartic/quartic45.json} and
\texttt{quartic/quartic45\_expanded.json}.

\subsection{Why the unguarded transport quotient is harmless}

The old fifty-variable system had independent ternary guards for
\(J,F,\alpha,z,\eta,\zeta\).  Delete the first four guards and their
twelve private variables; retain the last two; add \(v_q\) and the two
new triples.  Thus the variable count is
\(50-12+7=45\), and the residual count is \(29-4+3=28\).
Because these are domain changes, an algebraic identity alone would not
prove the required existential equivalence.

At a zero of \eqref{eq:quartic45}, the dyadic guard gives \(J>0\),
and the common-sign guard gives \(F,\alpha,C>0\).  The retained
guards give \(\eta,\zeta>0\).  The only missing original positivity
condition is \(z>0\).  Its recovery must respect the proof order.

In the inherited signed Pell proof, the early use of \(z>0\) was
solely to infer \(C\ge0\) from transport.  We already have the stronger
conclusion \(C>0\), independently.  Hence the same packing bounds follow:
\[
 q(q^2-1)<R<q^4-2,\qquad R-2>3q+1.
\]
The square-scale exclusion gives \(w>0\).  The existing odd-scale
argument, with \(s=2v_s+1\), excludes negative \(s\); it uses the
Pell equations, these bounds, and the exposed congruences
\(V\equiv-c\pmod{|f|}\), \(V\equiv-R\pmod{|c|}\), without
another use of \(z\).  Positive-scale rank recovery and the signed
main-root argument then give
\[
 X=2^R.
\]
For a detailed dependency trace see the included
\texttt{ODD\_SCALE\_27\_PROOF.md}, Sections 1--4, and
\texttt{SIGNED\_DOMAIN\_29.md}, through its signed-gamma and
ordinary-strong arguments.  This explicitly supplies \(C>0\) in place
of the original preliminary transport inference; it does not invoke
the final signed-domain theorem before restoring its guard.

Since \(R>3q\), we have \(w=2^R/q>q>F\).  The retained transport
equation now reads
\[
 (K_0+w)C=(z-1)(q-1)+F.
\]
Its left side exceeds \(F\), so \(z>1\).  All original positivity
conditions are therefore restored.  Choosing their ternary representations
produces a zero of the old fifty-variable quartic, and its inherited
compiler theorem yields a finite ZFC refutation.

Conversely, canonical accepting witnesses have \(q=2^t\) with odd
\(t\), and \(F,\alpha,C>0\).  Set \(v_q=2^{(t-1)/2}\) and use
the positive representations of \(F\alpha\) and \(FC\).  Every new
equation holds.  Consequently the ordinary mathematical argument gives,
formalizably in ZFC,
\[
 \text{there is no integer zero of }Q_{45}
 \quad\Longleftrightarrow\quad \operatorname{Con}(\mathrm{ZFC}).
\]

The portable finite audit checks all 28 residual degrees, the 45 live
coordinates, the exact expansion, \([\tau^4]Q_{45}=1\), the complete
fixed payloads, and the pinned inherited bytes.  The unbounded Pell and
compiler arguments remain explicit mathematical dependencies; this
section does not claim a new end-to-end Lean/Rocq certificate or an
exported object-language ZFC derivation.
